\chapter{Negative-Pressure Wound Therapy}

\section*{Introduction \& Clinical Indications}
Negative-pressure wound therapy (NPWT) applies controlled suction through a sealed dressing to remove exudate, reduce edema, support wound-edge contraction, and promote granulation. It can be used for selected acute traumatic wounds, dehisced incisions, chronic wounds, graft bolsters, and temporary management of an open abdomen. It is a wound-management adjunct, not a substitute for debridement, drainage, antibiotics when indicated, or definitive closure.

The wound must be assessed for perfusion, infection, necrosis, exposed structures, and the need for urgent source control. NPWT may be inappropriate over untreated osteomyelitis, an unrecognized enteric or urinary fistula, necrotic tissue, malignancy in the wound, or exposed vessels, nerves, anastomoses, or organs without protective coverage. Pressure settings and dressing intervals are individualized.

\section*{Relevant Surgical Anatomy}
Foam or gauze must contact the intended wound surface without being placed directly on delicate vessels, bowel, nerves, or tendon unless an appropriate nonadherent interface protects them. Tunneling and undermining should be documented because retained foam can cause persistent infection. The seal must avoid pressure injury at bony prominences and must not compromise adjacent skin or stomas.

\section*{Preoperative Workup \& Preparation}
Evaluate circulation, sensation, wound dimensions, tissue viability, exudate, odor, pain, nutrition, diabetes, anticoagulation, and infection. Debride devitalized tissue and drain collections before application. Discuss discomfort, bleeding, skin injury, odor, alarm failure, fluid loss, and the possibility that treatment will need to stop or change if the wound deteriorates. Confirm equipment, power or battery supply, and a plan for transfer and emergency dressing removal.

\section*{Surgical Technique \& Procedure}
The wound is cleansed and measured. A correctly shaped foam or gauze interface is placed, the surrounding skin is protected, and an occlusive film creates the seal. The tubing is connected to the pump and the prescribed pressure and mode are selected. The dressing should collapse appropriately and alarms should be tested. At changes, inspect the wound, count all materials, assess granulation and necrosis, and remeasure before deciding whether to continue, bridge, graft, or close.

\section*{Complications \& Risks}
Bleeding is the most serious risk, particularly over vascular structures or in anticoagulated patients. Other complications include pain, maceration, dermatitis, retained foam, loss of seal, infection, pressure injury, dehydration, and fistula formation. Sudden bright-red drainage, hemodynamic change, severe pain, fever, or persistent alarms requires immediate evaluation rather than simply increasing suction.

\section*{Postoperative Care \& Recovery}
Monitor output, seal, pain, skin, wound appearance, temperature, and systemic status. Dressing changes may require analgesia or anesthesia. Maintain nutrition and glycemic control, treat infection based on clinical assessment and cultures when appropriate, and involve surgery, wound specialists, nursing, and rehabilitation. Do not allow a disconnected or nonfunctioning system to remain unattended for prolonged periods.

\section*{Outcomes \& Prognosis}
NPWT can prepare selected wounds for delayed closure, grafting, or secondary healing and may reduce handling of exudate. Success depends on correct indication, source control, protection of vital structures, and timely transition to definitive treatment. Evidence and benefit vary by wound type; treatment should be reassessed at each change.

\section*{References}
\begin{enumerate}
\item International Wound Infection Institute. Wound infection guidance.
\item National Institute for Health and Care Excellence. Negative pressure wound therapy for the open abdomen and surgical wounds.
\item European Wound Management Association. Position documents on negative pressure wound therapy.
\end{enumerate}
